\section{Analyse fonctionnelle}

\subsection{Expression du besoin}

\subsubsection{Bête à cornes}

\begin{longtable}{|p{5cm}|p{8cm}|}
\hline
\textbf{Question} & \textbf{Réponse} \\
\hline
À qui le produit rend-il service ? &
À un utilisateur qui souhaite identifier une musique depuis un objet autonome. \\
\hline
Sur quoi agit-il ? &
Sur un extrait sonore capté dans l'environnement. \\
\hline
Dans quel but ? &
Reconnaître le morceau et afficher les informations principales. \\
\hline
\end{longtable}

Formulation du besoin :

\begin{quote}
Permettre à un utilisateur d'identifier une musique diffusée autour de lui à
l'aide d'un objet physique autonome, lisible et documenté.
\end{quote}

\subsection{Environnement du système}

Le système interagit avec les éléments extérieurs suivants :

\begin{itemize}
    \item utilisateur ;
    \item musique ambiante ;
    \item réseau Wi-Fi local ;
    \item service de reconnaissance musicale AudD ;
    \item écran e-paper ;
    \item boutons physiques ;
    \item batterie 2000 mAh et UPS HAT (C) ;
    \item navigateur web ;
    \item boîtier du prototype.
\end{itemize}

\subsection{Diagramme pieuvre textuel}

\begin{longtable}{|p{1.8cm}|p{8.8cm}|p{2.5cm}|}
\hline
\textbf{Code} & \textbf{Relation avec l'environnement} & \textbf{Type} \\
\hline
FP1 & Identifier une musique diffusée dans l'environnement. & Principale \\
\hline
FC1 & Être commandé simplement par l'utilisateur. & Contrainte \\
\hline
FC2 & Afficher une information lisible sur écran e-paper. & Contrainte \\
\hline
FC3 & Fonctionner sur Raspberry Pi avec ressources limitées. & Contrainte \\
\hline
FC4 & Communiquer avec un réseau Wi-Fi local. & Contrainte \\
\hline
FC5 & Utiliser des composants accessibles. & Contrainte \\
\hline
FC6 & Être maintenable et documenté. & Contrainte \\
\hline
FC7 & S'intégrer dans un objet au style cassette futuriste. & Contrainte \\
\hline
FC8 & Gérer les échecs de reconnaissance sans blocage. & Contrainte \\
\hline
\end{longtable}

\subsection{Fonctions de service}

\begin{longtable}{|p{2cm}|p{5cm}|p{6cm}|}
\hline
\textbf{Code} & \textbf{Fonction} & \textbf{Description} \\
\hline
FS1 & Capturer un extrait audio &
Le système enregistre un extrait court au format WAV. \\
\hline
FS2 & Reconnaître le morceau &
Le fichier audio est transmis à l'API AudD pour obtenir un résultat. \\
\hline
FS3 & Normaliser les métadonnées &
Le résultat est transformé en structure utilisable par l'affichage. \\
\hline
FS4 & Afficher l'état ou le résultat &
L'écran e-paper affiche l'état courant ou le morceau reconnu. \\
\hline
FS5 & Configurer le Wi-Fi &
Une interface Flask permet de choisir un réseau et de saisir le mot de passe. \\
\hline
FS6 & Réagir aux boutons &
Les boutons physiques déclenchent le changement de mode et la détection. \\
\hline
\end{longtable}

\subsection{Scénarios d'utilisation}

\subsubsection{Scénario nominal}

\begin{enumerate}
    \item L'utilisateur allume S.Y.S.T.U.S.
    \item Le système vérifie l'état du Wi-Fi.
    \item Si le Wi-Fi est disponible, le mode \texttt{running} est affiché.
    \item L'utilisateur appuie sur le bouton d'action.
    \item Le système enregistre un extrait audio. L'écoute n'est pas automatique en continu : chaque écoute est déclenchée par un clic bouton.
    \item L'extrait est envoyé au service AudD.
    \item Les métadonnées retournées sont filtrées.
    \item L'écran e-paper affiche le titre, l'artiste, l'album, l'année et un QR code de lien si disponible.
\end{enumerate}

\subsubsection{Scénario configuration Wi-Fi}

\begin{enumerate}
    \item Le système démarre sans connexion Wi-Fi active.
    \item Le mode \texttt{setup} est sélectionné.
    \item L'écran affiche un QR code vers l'interface locale.
    \item L'utilisateur ouvre l'interface web.
    \item L'application scanne les réseaux disponibles avec \texttt{nmcli}.
    \item L'utilisateur saisit le SSID et le mot de passe.
    \item Le système tente la connexion et revient au mode \texttt{running} si elle réussit.
\end{enumerate}

Dans l'implémentation actuelle, le bouton de mode permet aussi de basculer
manuellement vers le mode de mise en place du Wi-Fi.

\subsubsection{Scénario d'échec de reconnaissance}

\begin{enumerate}
    \item L'utilisateur lance une écoute.
    \item L'audio est capturé, mais aucun morceau n'est reconnu.
    \item Le système affiche un état \textit{No result}.
    \item L'utilisateur peut relancer une écoute.
\end{enumerate}

\subsection{Chaîne fonctionnelle}

\begin{lstlisting}
Utilisateur
  -> Bouton d'action
  -> DetectionController
  -> AudioRecorder
  -> arecord / fichier WAV temporaire
  -> AudDClient
  -> Metadonnees normalisees
  -> Module display
  -> Ecran e-paper
\end{lstlisting}

\subsection{Fonctions critiques}

\begin{longtable}{|p{4cm}|p{3cm}|p{6cm}|}
\hline
\textbf{Fonction} & \textbf{Criticité} & \textbf{Justification} \\
\hline
Capture audio & Critique & Sans signal exploitable, aucune reconnaissance n'est possible. \\
\hline
Reconnaissance musicale & Critique & C'est le coeur de la promesse utilisateur. \\
\hline
Affichage e-paper & Critique & C'est la sortie principale du prototype. \\
\hline
Gestion Wi-Fi & Importante & Nécessaire pour l'API et la configuration. \\
\hline
Boutons physiques & Importante & Rend l'objet manipulable sans ordinateur. \\
\hline
Boîtier & Moyenne & Important pour l'expérience, mais non bloquant pour valider le flux. \\
\hline
\end{longtable}
